\section{边疆、经学与宫廷：守成朝廷的多重成本}
\label{sec:eastern-han-frontier-learning-court}

光武之后的东汉并非只是在雒阳守成。永平八年，楚王刘英因被告谋反获减死并被徙，诏令提及黄老与佛教供养。这是汉地佛教进入官方文字的一条早期线索，却不能由一纸诏书推断全国已经形成统一的佛教制度，术语的文本层次和实际宗教实践仍有距离。事件同时说明，宗室、礼制与新出现的供养语言都已能够成为朝廷判断忠诚的对象。

\subsection{西域不是一条可以永久占有的道路}

七三年窦固等出击北匈奴、取伊吾卢，班超奉使鄯善，东汉重新进入西域；两年后北匈奴攻车师，西域都护陈睦遇害，朝廷一度撤出据点。两件事放在一起，比一条“经营西域”的总述更能显示问题所在：绿洲之间的使节、驻军和地方王权需要持续的道路、粮秣与联盟，名义上的都护并不自动变成直接统治。各路战果、撤军范围和车师、北匈奴各部在危机中的位置都有异说，不能把前进和后退画成边界清楚的地图。

章帝朝以后，西羌在陇西、金城等地的反叛又持续增加边郡军费和迁徙压力。史书中的“羌”首先是当时官府的分类，不能被想作内部完全一致的社会实体。对于雒阳而言，西域与陇右共同提出同一个难题：远方驻军、地方首领与道路安全会不断要求资源，而中央所能获得的户口、粮食和官员并非没有边界。

\subsection{经学的会议与外戚的军功}

七九年章帝召儒者于白虎观讨论五经异同。会议存在可靠，现存《白虎通》与当时讨论记录之间的编定关系却不可忽略。它表现了朝廷希望以经典语言协调礼制、名分和官员资格；经学进入官方话语，并不说明不同师说已经消失，更不说明地方教育被同一套课程覆盖。

八八年和帝即位，窦太后临朝，窦氏成为人事和军事任命的中心。次年窦宪北出稽落山，重创北匈奴并以燕然山铭宣示战功。战役年月与铭文存在可相互校核，兵数、斩获和铭文的功绩措辞则正是需要警惕的部分：铭文是公开的胜利表达，不是独立的战场统计。军功抬升了窦氏的政治位置，也使幼帝朝廷更难区分外戚的边防贡献与其对中枢的支配。

\subsection{班超的网络与宫廷的反转}

九一年班超受命为西域都护，东汉在部分绿洲重新建立使节和军事网络；九七年甘英西使，抵安息属地而未能到达大秦。使节的发生和所见的地理知识可由《后汉书》相关传记、〈西域传〉相互参照，但到达地点、海路信息和“大秦”的识别皆有争议。更不能由这一行程推定汉与罗马已经建立直接外交。它所显出的，是汉廷经由中介政权取得远方信息的能力，以及这种能力依赖商路与地方合作的限度。

九二年和帝与宦官郑众诛除窦宪集团。政变事实可靠，和帝的自主程度和郑众的角色仍须保留疑问；其后果是宦官因参与清除外戚而进入更关键的政治位置。班超于一〇二年返洛阳后去世，西域经营失去长期协调者。个人网络不能完全等同于制度，后任也并非没有行动；不过，朝廷在边疆和宫廷两端都越来越依靠能够替中枢临时解决问题的个人与集团。

\subsection{来源的不同亮度}

《后汉书》帝纪、楚王英传、窦固传、班超传、西羌传和〈西域传〉提供主干，《后汉纪》帮助校看年序；《牟子理惑论》传本、《白虎通》、燕然山铭和相关考古材料分别照出宗教、经学、军功与空间的不同层面。它们的成书距离、文本层次和宣示目的并不相同。本节因此把宗室案件、白虎观、边疆驻军、使节和外戚政变写为守成国家的相互牵制，而不把任何一类材料扩展成全国社会的直接证明。
